\chapter{Anal Fissure}

\section*{Definition and Overview}
An anal fissure is a small tear or split in the lining of the anal canal, usually at the point where it meets the surrounding skin. It is a common and very painful condition, typically felt as a sharp, burning pain when passing stool, followed by a throbbing ache that can last for hours.

Most fissures are caused by passing a hard or large stool, and the majority heal with simple measures such as softening the stool and relaxing the muscle around the anus. A fissure is called chronic when it lasts longer than about six to eight weeks despite treatment.

\section*{Epidemiology}
Anal fissures are very common and affect people of all ages, including children and older adults. They occur about equally in men and women, although in women they are more frequent during and after pregnancy and childbirth. Roughly one in ten people will experience an anal fissure at some point in their lives.

Most fissures are acute and heal quickly, but a significant minority become chronic. Fissures are more common in people with constipation, hard stools, frequent loose stools, or a tendency to strain on the toilet.

\section*{Aetiology and Pathophysiology}
An anal fissure usually begins as a mechanical tear in the lining of the anal canal:
\begin{itemize}
    \item \textbf{Passing a hard or large stool:} the most common cause, often with straining
    \item \textbf{Prolonged diarrhoea:} frequent, urgent, loose stools can injure the lining
    \item \textbf{Childbirth:} a tear may occur during delivery
    \item \textbf{Tight anal muscle:} the sphincter goes into spasm, pulling the wound edges apart and reducing the blood supply, which slows healing
    \item \textbf{Anal trauma:} from anal sex, constipation, or insertion of objects
    \item \textbf{Underlying conditions:} inflammatory bowel disease (especially Crohn's disease), and less commonly infection or cancer, can cause fissures or prevent them from healing
\end{itemize}
The cycle of tearing, muscle spasm, and reduced blood flow is the reason some fissures fail to heal and become chronic.

\section*{Clinical Features}
\begin{itemize}
    \item Sharp, tearing pain during or just after passing stool
    \item A burning or throbbing ache that may last for minutes to hours afterwards
    \item Bright red blood on the toilet paper or on the surface of the stool
    \item A visible split or tear at the anal opening, often at the back (posterior midline)
    \item A skin tag (sentinel tag) at the edge of a chronic fissure
    \item Fear of passing stool, sometimes leading to withholding and worsening constipation, especially in children
\end{itemize}

\section*{Differential Diagnosis}
Pain and bleeding around the anus can be caused by other conditions:
\begin{itemize}
    \item Haemorrhoids (piles)
    \item Perianal abscess or fistula
    \item Anal cancer or pre-cancerous changes, which can resemble a non-healing fissure
    \item Inflammatory bowel disease, including Crohn's disease, which causes atypical or multiple fissures
    \item Sexually transmitted infections, such as herpes, and other skin conditions
    \item Anal trauma from other causes
\end{itemize}

\section*{When to Seek Medical Care}
See a doctor if:
\begin{itemize}
    \item Pain or bleeding from the back passage persists for more than a few days
    \item The pain is severe or not helped by simple measures
    \item The fissure has not healed after several weeks of self-care
    \item You have an unusual fissure, such as one at the side of the anus, several fissures, or one that is painless
    \item You have symptoms suggesting Crohn's disease, such as chronic diarrhoea, weight loss, or abdominal pain
\end{itemize}

\textbf{Call 000 (or local emergency services) for:}
\begin{itemize}
    \item Heavy bleeding from the back passage
    \item Sudden severe pain or collapse
    \item An inability to pass urine or stool with worsening abdominal pain
\end{itemize}

\section*{Diagnosis and Investigations}
\begin{itemize}
    \item \textbf{History:} description of the pain, bleeding, and bowel habit, including whether this has happened before
    \item \textbf{Examination:} gentle inspection of the anal area, which is often all that is needed to make the diagnosis
    \item \textbf{Digital rectal examination:} performed if the fissure is not too painful, to assess the sphincter muscle and surrounding tissue
    \item \textbf{Proctoscopy:} direct viewing of the anal canal, sometimes needed to exclude other causes
    \item \textbf{Examination under anaesthetic:} used for chronic or suspicious fissures when examination is too painful
    \item \textbf{Biopsy:} taken from a chronic or atypical fissure if there is concern about another diagnosis
    \item \textbf{Colonoscopy:} may be arranged if inflammatory bowel disease is suspected
\end{itemize}

\section*{Management}
\begin{itemize}
    \item \textbf{Stool softening:} more fluid, dietary fibre, and stool softeners such as psyllium so that passing stool is less painful
    \item \textbf{Pain relief:} warm baths (sitz baths) several times a day and simple painkillers as advised
    \item \textbf{Skin care:} keeping the area clean and dry with gentle wiping
    \item \textbf{Sphincter-relaxing creams and ointments:} nitroglycerin or calcium-channel-blocker creams applied around the anus to relax the muscle and improve blood flow
    \item \textbf{Botulinum toxin (Botox) injection:} into the sphincter to relax it in chronic cases
    \item \textbf{Surgery:} lateral internal sphincterotomy, a small operation to divide part of the sphincter muscle, for fissures that do not heal with other treatment
    \item \textbf{Treating the underlying cause:} such as managing Crohn's disease or an infection
    \item \textbf{Follow-up:} review to confirm healing and address any persisting symptoms
\end{itemize}

\section*{Complications}
\begin{itemize}
    \item A chronic fissure that fails to heal despite treatment
    \item Recurrent fissures with repeated bowel trauma
    \item A sentinel skin tag, which may be mistaken for a haemorrhoid
    \item Painful spasm of the anal muscle, which prolongs the problem
    \item Scarring or narrowing (stricture) of the anal canal, rarely, after chronic fissures or surgery
    \item A small risk of incontinence of gas or stool after sphincter surgery
    \item Masking of an underlying condition such as inflammatory bowel disease or, rarely, cancer
\end{itemize}

\section*{Prognosis and Outcomes}
Most anal fissures heal with simple conservative measures. Roughly half or more heal with stool softening and warm baths alone, and the majority that remain heal with the addition of sphincter-relaxing creams or injections. Surgery has a very high success rate for chronic fissures, although it carries a small risk of affecting bowel control. Recurrence is possible, especially if constipation or straining continues, but is less likely when the person maintains soft stools and good anal care.

\section*{Prevention and Self-Care}
\begin{itemize}
    \item Keep stools soft with plenty of water and a diet rich in fibre, including fruits, vegetables, and whole grains
    \item Do not delay or strain when you need to open your bowels
    \item Sit on the toilet only for as long as needed and avoid pushing
    \item Use stool softeners in the short term if you are prone to hard stools
    \item Take warm baths to soothe the area and relax the anal muscle
    \item Keep the anal area clean and dry, wiping gently and avoiding perfumed products
    \item Have persistent symptoms reviewed by a doctor, especially with recurrent fissures or a family history of bowel disease
    \item Women should be aware of the higher risk of fissures around childbirth and seek early advice if symptoms develop
\end{itemize}

\section*{Sources and Further Reading}
The following resources provide reliable, peer-reviewed and patient-orientated information on anal fissures.

\begin{itemize}
    \item NHS. \textit{Overview of anal fissure, symptoms, and treatment.} https://www.nhs.uk/conditions/
    \item Mayo Clinic. \textit{Guide to the diagnosis and management of anal fissures.} https://www.mayoclinic.org/
    \item MSD Manual. \textit{Professional and patient information on anal fissures.} https://www.msdmanuals.com/
    \item MedlinePlus. \textit{Health information on anal fissure.} https://medlineplus.gov/
    \item NICE. \textit{Clinical guidance on the management of anal fissure.} https://www.nice.org.uk/
\end{itemize}
